%% Abstract
Diffusion and flow-matching transformers generate high-quality images and videos and are increasingly used in engineering design, but sampling them requires tens to thousands of sequential evaluations of a large network. Training-free feature caching reduces this cost by skipping the network on some steps; every cache must decide \emph{when} to compute and \emph{what} to write on the skipped steps. This project develops two caches that improve both decisions.

\emph{ReSPect} (Residual Spectral Prediction) couples a spectrally aligned gate, derived from the linear minimum-mean-square-error (Wiener) denoiser, with a forecaster of the transformer \emph{residual} that blends a ridge-regularised Chebyshev fit of the residual history with a first-order discrete Taylor correction. We show that its forecast error admits a bound independent of the skip length, whereas the bounds of copy-reuse and Taylor extrapolation grow with it, and we derive the blend weight and warm-up analytically. ReSPect achieves speedups of $3.30\times$, $2.63\times$ and $2.77\times$ on FLUX.1-dev, Wan2.1-1.3B and HunyuanVideo with lower output error than existing caching baselines at matched or lower compute, improving PSNR by up to 3.46, 11.4 and 9.9\,dB.

\emph{KalSea} (Kalman-Corrected Spectral-Evolution-Aware Caching) observes that all existing fillers are open loop: the error made on skipped steps is integrated into the sample and never corrected. Behind the same gate, a damped-trend Kalman filter forecasts the residual, and \emph{predictor--corrector gap smoothing} (PCGS) uses the Rauch--Tung--Striebel smoother to correct every skipped gap once the next step is computed, re-integrating it through the output layers only, with no additional transformer evaluations. On FLUX.1-dev, KalSea reaches 27.72\,dB PSNR with respect to the 50-step sampler at 18.1 of 50 evaluations (SeaCache: 26.29\,dB at 18.4) and 22.53 vs.\ 21.33\,dB at 13 evaluations; PCGS alone contributes 0.81 and 0.61\,dB, and KalSea improves over SeaCache on HunyuanVideo and Wan2.1 by up to 3.2 and 2.3\,dB in the same setup.

Finally, we carry both caches over to two open-weight diffusion models from mechanical engineering, Optimize Any Topology for structural topology optimisation and DiffusionPDE for the reconstruction of a flow past a cylinder, by generalising the residual split to U-Net denoisers and the corrector to DDIM with classifier-free guidance and to gradient-guided second-order samplers. On DiffusionPDE's Navier--Stokes forward and inverse problems, KalSea needs $1.7$--$2.6\times$ fewer network evaluations than the 2000-step sampler, recovers the flow with 3.3\% error on the forward problem of the first released case, and is more accurate than ReSPect in all eight comparisons, showing that the closed loop is essential when guidance gradients are approximated on skipped steps.
